\documentclass[11pt,a4paper]{article}

%========================
% PACKAGES
%========================
\usepackage[utf8]{inputenc}
\usepackage[T1]{fontenc}
\usepackage[french]{babel}
\usepackage[a4paper,margin=2cm]{geometry}
\usepackage{amsmath,amssymb,amsthm}
\usepackage{lmodern}
\usepackage{fancyhdr}
\usepackage{lastpage}
\usepackage{enumitem}
\usepackage{array,tabularx}
\usepackage{tikz}
\usepackage{pgfplots}
\usepackage[most]{tcolorbox}
\usepackage{xcolor}
\usepackage{hyperref}

\pgfplotsset{compat=1.18}

%========================
% COULEURS
%========================
\definecolor{bleu}{RGB}{30,90,170}
\definecolor{vert}{RGB}{30,130,80}
\definecolor{rouge}{RGB}{180,50,50}
\definecolor{orange}{RGB}{220,130,30}
\definecolor{gris}{RGB}{245,245,245}

%========================
% MISE EN PAGE
%========================
\setlength{\headheight}{14pt}
\pagestyle{fancy}
\fancyhf{}
\fancyhead[L]{Première spécialité mathématiques}
\fancyhead[C]{Fonction exponentielle}
\fancyhead[R]{Page \thepage/\pageref{LastPage}}
\fancyfoot[C]{Cours complet -- Fonction exponentielle}

\hypersetup{
    colorlinks=true,
    linkcolor=bleu,
    urlcolor=bleu
}

%========================
% ENVIRONNEMENTS
%========================
\tcbset{
  colback=gris,
  colframe=bleu,
  arc=2mm,
  boxrule=0.7pt,
  left=2mm,
  right=2mm,
  top=1mm,
  bottom=1mm
}

\newtcolorbox{definitionbox}{
  title=\textbf{Définition},
  colback=blue!4,
  colframe=bleu
}

\newtcolorbox{propriete}{
  title=\textbf{Propriété},
  colback=green!5,
  colframe=vert
}

\newtcolorbox{theoremebox}{
  title=\textbf{Théorème},
  colback=blue!5,
  colframe=bleu
}

\newtcolorbox{methode}{
  title=\textbf{Méthode},
  colback=orange!8,
  colframe=orange
}

\newtcolorbox{retenir}{
  title=\textbf{À retenir},
  colback=yellow!10,
  colframe=orange
}

\newtcolorbox{erreur}{
  title=\textbf{Erreur classique},
  colback=red!5,
  colframe=rouge
}

\newtcolorbox{pointbac}{
  title=\textbf{Point contrôle},
  colback=purple!5,
  colframe=purple!70!black
}

\newtcolorbox{exemplebox}{
  title=\textbf{Exemple corrigé},
  colback=white,
  colframe=black!60
}

\newtcolorbox{remarque}{
  title=\textbf{Remarque},
  colback=gray!10,
  colframe=black!50
}

\newcounter{exo}
\newenvironment{exercice}[1][]{
\refstepcounter{exo}
\par\medskip
\noindent\textbf{\textcolor{bleu}{Exercice \theexo.}} \textit{#1}
\par\smallskip
}{\medskip}

\newcounter{corr}
\newenvironment{correction}[1][]{
\refstepcounter{corr}
\par\medskip
\noindent\textbf{\textcolor{vert}{Correction de l'exercice \thecorr.}} \textit{#1}
\par\smallskip
}{\medskip}

%========================
% COMMANDES
%========================
\newcommand{\R}{\mathbb{R}}
\newcommand{\e}{\mathrm{e}}
\newcommand{\dd}{\mathrm{d}}

%========================
% DOCUMENT
%========================
\begin{document}

%========================
% PAGE DE TITRE
%========================
\begin{titlepage}
\centering
\vspace*{2cm}

{\Huge \textbf{Cours complet}}\\[0.4cm]
{\LARGE \textbf{Fonction exponentielle}}\\[0.4cm]
{\Large Première spécialité mathématiques}\\[1cm]

\begin{tcolorbox}[colback=blue!5,colframe=bleu,width=0.85\textwidth]
\centering
\Large
Définitions, propriétés, méthodes, exemples corrigés, exercices progressifs et corrections détaillées.
\end{tcolorbox}

\vfill

{\large Polycopié de cours}\\
{\large Niveau Première spécialité mathématiques}

\vfill
\end{titlepage}

\tableofcontents
\newpage

%================================================
\section{Introduction à la fonction exponentielle}

\subsection*{A. Objectifs du chapitre}

Dans ce chapitre, on introduit une nouvelle fonction fondamentale : la fonction exponentielle. Elle sert à modéliser des phénomènes où une quantité évolue proportionnellement à elle-même.

À la fin du chapitre, il faut savoir :
\begin{itemize}
    \item comprendre l'idée de croissance exponentielle ;
    \item utiliser les propriétés de calcul de l'exponentielle ;
    \item dériver des fonctions simples contenant une exponentielle ;
    \item résoudre des équations du type $\e^A=\e^B$ ;
    \item étudier les variations d'une fonction contenant $\e^x$ ;
    \item interpréter des modèles simples de croissance ou de décroissance.
\end{itemize}

\subsection*{B. Croissance additive et croissance multiplicative}

\begin{definitionbox}
Une \textbf{croissance additive} consiste à ajouter régulièrement une même quantité.

Par exemple, si une population augmente de $100$ individus chaque année, l'évolution est additive.
\end{definitionbox}

\begin{definitionbox}
Une \textbf{croissance multiplicative} consiste à multiplier régulièrement par un même facteur.

Par exemple, si une population augmente de $5\%$ chaque année, elle est multipliée chaque année par $1,05$.
\end{definitionbox}

\begin{exemplebox}
Un capital de $1000$ euros augmente de $4\%$ chaque année.

Après un an :
\[
1000 \times 1,04.
\]

Après deux ans :
\[
1000 \times 1,04^2.
\]

Après $n$ années :
\[
1000 \times 1,04^n.
\]

On retrouve ici une suite géométrique.
\end{exemplebox}

\subsection*{C. Lien avec les suites géométriques}

\begin{propriete}
Une suite géométrique de premier terme $u_0$ et de raison $q$ est définie par :
\[
u_n=u_0q^n.
\]

Elle modélise une évolution multiplicative discrète, c'est-à-dire étape par étape.
\end{propriete}

La fonction exponentielle prolonge cette idée à des valeurs réelles du temps. Elle permet de modéliser des évolutions continues.

\begin{retenir}
Les suites géométriques modélisent une évolution multiplicative discrète.

La fonction exponentielle modélise une évolution multiplicative continue.
\end{retenir}

\subsection*{D. Exemples de situations exponentielles}

On rencontre l'exponentielle dans de nombreux domaines :
\begin{itemize}
    \item intérêts composés ;
    \item évolution d'une population ;
    \item décroissance radioactive ;
    \item dilution d'une substance ;
    \item refroidissement d'un corps ;
    \item propagation d'une information ;
    \item modèles économiques ou biologiques.
\end{itemize}

\begin{erreur}
Une croissance exponentielle n'est pas simplement une croissance rapide.

Ce qui la caractérise, c'est que l'évolution est proportionnelle à la quantité présente.
\end{erreur}

%================================================
\section{Définition de la fonction exponentielle}

\subsection*{A. Définition}

\begin{definitionbox}
On appelle \textbf{fonction exponentielle} l'unique fonction dérivable sur $\R$ telle que :
\[
\exp'(x)=\exp(x)
\]
pour tout réel $x$, et :
\[
\exp(0)=1.
\]
\end{definitionbox}

Autrement dit, la fonction exponentielle est une fonction égale à sa propre dérivée.

\begin{remarque}
L'existence et l'unicité d'une telle fonction sont admises au niveau Première.
\end{remarque}

\subsection*{B. Notation $\e^x$}

\begin{definitionbox}
On note :
\[
\e=\exp(1).
\]

Pour tout réel $x$, on écrit :
\[
\exp(x)=\e^x.
\]
\end{definitionbox}

Ainsi :
\[
\exp(0)=\e^0=1,
\qquad
\exp(1)=\e^1=\e.
\]

Le nombre $\e$ vaut approximativement :
\[
\e \approx 2,718.
\]

\subsection*{C. Interprétation graphique}

La condition :
\[
\exp'(x)=\exp(x)
\]
signifie que le coefficient directeur de la tangente à la courbe au point d'abscisse $x$ est égal à l'ordonnée de ce point.

En particulier, au point d'abscisse $0$ :
\[
\exp(0)=1
\]
et
\[
\exp'(0)=1.
\]

Donc la tangente à la courbe en $0$ a pour équation :
\[
y=x+1.
\]

\begin{center}
\begin{tikzpicture}
\begin{axis}[
    axis lines=middle,
    xmin=-3,xmax=3,
    ymin=-0.5,ymax=6,
    grid=both,
    width=11cm,
    height=7cm,
    xlabel={$x$},
    ylabel={$y$},
    samples=200
]
\addplot[blue,thick,domain=-3:2] {exp(x)};
\addlegendentry{$y=\e^x$}
\addplot[red,dashed,domain=-2:3] {x+1};
\addlegendentry{$y=x+1$}
\addplot[only marks] coordinates {(0,1)};
\end{axis}
\end{tikzpicture}
\end{center}

%================================================
\section{Propriétés algébriques de l'exponentielle}

\subsection*{A. Relation fondamentale}

\begin{theoremebox}
Pour tous réels $a$ et $b$ :
\[
\e^{a+b}=\e^a\e^b.
\]
\end{theoremebox}

Cette relation est la propriété fondamentale de l'exponentielle.

\subsection*{B. Conséquences}

\begin{propriete}
Pour tous réels $a$ et $b$ :
\[
\e^{-a}=\frac{1}{\e^a},
\]
\[
\e^{a-b}=\frac{\e^a}{\e^b},
\]
et pour tout entier relatif $n$ :
\[
\e^{na}=(\e^a)^n.
\]
\end{propriete}

\subsection*{C. Démonstrations accessibles}

\begin{proof}
On sait que :
\[
\e^{a+(-a)}=\e^a\e^{-a}.
\]
Or :
\[
a+(-a)=0,
\]
donc :
\[
\e^0=\e^a\e^{-a}.
\]
Comme $\e^0=1$, on obtient :
\[
1=\e^a\e^{-a}.
\]
Comme $\e^a \neq 0$, on peut diviser par $\e^a$ :
\[
\e^{-a}=\frac{1}{\e^a}.
\]
\end{proof}

Pour la propriété du quotient :
\[
\e^{a-b}=\e^{a+(-b)}=\e^a\e^{-b}.
\]
Or :
\[
\e^{-b}=\frac{1}{\e^b}.
\]
Donc :
\[
\e^{a-b}=\frac{\e^a}{\e^b}.
\]

\subsection*{D. Méthode de simplification}

\begin{methode}
Pour simplifier une expression contenant des exponentielles :
\begin{enumerate}
    \item Repérer les produits et les quotients d'exponentielles.
    \item Utiliser :
    \[
    \e^a\e^b=\e^{a+b}.
    \]
    \item Utiliser :
    \[
    \frac{\e^a}{\e^b}=\e^{a-b}.
    \]
    \item Réduire l'exposant.
\end{enumerate}
\end{methode}

\begin{exemplebox}
Simplifier :
\[
A=\frac{\e^{2x+1}\e^{3x-4}}{\e^{x+2}}.
\]

On regroupe les exponentielles du numérateur :
\[
\e^{2x+1}\e^{3x-4}=\e^{(2x+1)+(3x-4)}=\e^{5x-3}.
\]

Donc :
\[
A=\frac{\e^{5x-3}}{\e^{x+2}}.
\]

On utilise la propriété du quotient :
\[
A=\e^{(5x-3)-(x+2)}.
\]

Ainsi :
\[
A=\e^{4x-5}.
\]
\end{exemplebox}

\begin{erreur}
Il ne faut pas écrire :
\[
\e^{a+b}=\e^a+\e^b.
\]

La bonne formule est :
\[
\e^{a+b}=\e^a\e^b.
\]
\end{erreur}

%================================================
\section{Signe, dérivée et variations de la fonction exponentielle}

\subsection*{A. Positivité}

\begin{theoremebox}
Pour tout réel $x$ :
\[
\e^x>0.
\]
\end{theoremebox}

\begin{remarque}
L'exponentielle ne s'annule jamais. La courbe de la fonction exponentielle reste toujours au-dessus de l'axe des abscisses.
\end{remarque}

\subsection*{B. Dérivée}

\begin{theoremebox}
Pour tout réel $x$ :
\[
(\e^x)'=\e^x.
\]
\end{theoremebox}

C'est la propriété caractéristique de la fonction exponentielle.

\subsection*{C. Variations}

Comme :
\[
\e^x>0
\]
pour tout réel $x$, on a :
\[
(\e^x)'>0.
\]

Donc la fonction exponentielle est strictement croissante sur $\R$.

\begin{center}
\begin{tabular}{c|ccc}
$x$ & $-\infty$ & & $+\infty$\\
\hline
$(\e^x)'$ & & $+$ & \\
\hline
$\e^x$ & & $\nearrow$ &
\end{tabular}
\end{center}

\begin{propriete}
Pour tous réels $a$ et $b$ :
\[
\e^a=\e^b \Longleftrightarrow a=b.
\]

De plus :
\[
\e^a<\e^b \Longleftrightarrow a<b.
\]
\end{propriete}

\subsection*{D. Exemple}

\begin{exemplebox}
Résoudre :
\[
\e^{2x-3}=\e^{5}.
\]

Comme la fonction exponentielle est strictement croissante, elle est injective. Donc :
\[
2x-3=5.
\]

Ainsi :
\[
2x=8,
\]
d'où :
\[
x=4.
\]

La solution est :
\[
S=\{4\}.
\]
\end{exemplebox}

\begin{exemplebox}
Résoudre :
\[
\e^{3x+1}<\e^{7}.
\]

Comme la fonction exponentielle est strictement croissante :
\[
3x+1<7.
\]

Donc :
\[
3x<6,
\]
et :
\[
x<2.
\]

La solution est :
\[
S=]-\infty;2[.
\]
\end{exemplebox}

\begin{erreur}
On ne résout pas une équation du type :
\[
\e^x=2
\]
avec le logarithme si le logarithme n'a pas encore été étudié.

En Première, on privilégie les équations de la forme :
\[
\e^A=\e^B.
\]
\end{erreur}

%================================================
\section{Représentation graphique de la fonction exponentielle}

\subsection*{A. Points remarquables}

La courbe de la fonction exponentielle passe par :
\[
(0;1),
\qquad
(1;\e),
\qquad
(-1;\frac{1}{\e}).
\]

\subsection*{B. Allure générale}

\begin{center}
\begin{tikzpicture}
\begin{axis}[
    axis lines=middle,
    xmin=-4,xmax=4,
    ymin=-0.5,ymax=8,
    grid=both,
    width=12cm,
    height=7cm,
    xlabel={$x$},
    ylabel={$y$},
    samples=200
]
\addplot[blue,thick,domain=-4:2.1] {exp(x)};
\addlegendentry{$y=\e^x$}
\addplot[only marks] coordinates {(0,1) (1,2.718) (-1,0.367)};
\end{axis}
\end{tikzpicture}
\end{center}

\begin{retenir}
La fonction exponentielle :
\begin{itemize}
    \item est définie sur $\R$ ;
    \item est strictement positive ;
    \item est strictement croissante ;
    \item passe par le point $(0;1)$ ;
    \item a pour tangente en $0$ la droite d'équation $y=x+1$.
\end{itemize}
\end{retenir}

\subsection*{C. Lecture graphique}

Puisque la fonction exponentielle est strictement croissante :
\[
\e^x=1 \Longleftrightarrow x=0.
\]

De même :
\[
\e^x>\e^2 \Longleftrightarrow x>2.
\]

\begin{pointbac}
Dans les exercices, la stricte croissance de l'exponentielle permet souvent de transformer une inéquation exponentielle en inéquation classique.
\end{pointbac}

%================================================
\section{Dérivation de fonctions simples avec exponentielle}

\subsection*{A. Fonctions du type $ae^x+b$}

Soit :
\[
f(x)=a\e^x+b.
\]

Alors :
\[
f'(x)=a\e^x.
\]

\begin{exemplebox}
Soit :
\[
f(x)=3\e^x-5.
\]

Alors :
\[
f'(x)=3\e^x.
\]

Comme $\e^x>0$, on a :
\[
f'(x)>0.
\]

Donc $f$ est strictement croissante sur $\R$.
\end{exemplebox}

\subsection*{B. Fonctions du type $\e^x+mx+p$}

Soit :
\[
f(x)=\e^x+mx+p.
\]

Alors :
\[
f'(x)=\e^x+m.
\]

L'étude du signe dépend de la valeur de $m$.

\begin{exemplebox}
Soit :
\[
f(x)=\e^x-x.
\]

On dérive :
\[
f'(x)=\e^x-1.
\]

Or :
\[
\e^x>1 \Longleftrightarrow x>0,
\]
\[
\e^x=1 \Longleftrightarrow x=0,
\]
\[
\e^x<1 \Longleftrightarrow x<0.
\]

Donc :
\[
f'(x)<0 \text{ sur } ]-\infty;0[,
\]
\[
f'(0)=0,
\]
\[
f'(x)>0 \text{ sur } ]0;+\infty[.
\]

Ainsi, $f$ est décroissante sur $]-\infty;0]$ puis croissante sur $[0;+\infty[$.

De plus :
\[
f(0)=\e^0-0=1.
\]

La fonction admet donc un minimum égal à $1$ en $x=0$.
\end{exemplebox}

\subsection*{C. Fonctions du type $P(x)e^x$}

En Première, on peut rencontrer des fonctions du type :
\[
f(x)=P(x)\e^x,
\]
où $P$ est un polynôme simple.

On utilise alors la formule du produit :
\[
(uv)'=u'v+uv'.
\]

\begin{exemplebox}
Soit :
\[
f(x)=x\e^x.
\]

On pose :
\[
u(x)=x,
\qquad
v(x)=\e^x.
\]

Alors :
\[
u'(x)=1,
\qquad
v'(x)=\e^x.
\]

Donc :
\[
f'(x)=1\times \e^x+x\times \e^x.
\]

Ainsi :
\[
f'(x)=(x+1)\e^x.
\]

Comme $\e^x>0$, le signe de $f'(x)$ est celui de $x+1$.
\end{exemplebox}

\begin{methode}
Pour étudier les variations d'une fonction contenant $\e^x$ :
\begin{enumerate}
    \item Calculer la dérivée.
    \item Factoriser si possible par $\e^x$.
    \item Utiliser le fait que $\e^x>0$.
    \item Étudier le signe du facteur restant.
    \item Construire le tableau de variations.
\end{enumerate}
\end{methode}

\begin{erreur}
Quand une dérivée est de la forme :
\[
f'(x)=(x-2)\e^x,
\]
il ne faut pas étudier séparément $\e^x$ comme si son signe pouvait changer.

On sait que :
\[
\e^x>0.
\]

Donc le signe de $f'(x)$ est uniquement celui de $x-2$.
\end{erreur}

%================================================
\section{Équations et inéquations avec exponentielle}

\subsection*{A. Équations du type $\e^A=\e^B$}

\begin{propriete}
Pour tous réels $A$ et $B$ :
\[
\e^A=\e^B \Longleftrightarrow A=B.
\]
\end{propriete}

\begin{exemplebox}
Résoudre :
\[
\e^{x^2-3x}=\e^{4}.
\]

On utilise l'injectivité de l'exponentielle :
\[
x^2-3x=4.
\]

Donc :
\[
x^2-3x-4=0.
\]

On factorise :
\[
x^2-3x-4=(x-4)(x+1).
\]

Ainsi :
\[
(x-4)(x+1)=0.
\]

Donc :
\[
x=4 \quad \text{ou} \quad x=-1.
\]

La solution est :
\[
S=\{-1;4\}.
\]
\end{exemplebox}

\subsection*{B. Inéquations du type $\e^A<\e^B$}

\begin{propriete}
Pour tous réels $A$ et $B$ :
\[
\e^A<\e^B \Longleftrightarrow A<B.
\]
\end{propriete}

\begin{exemplebox}
Résoudre :
\[
\e^{x^2}<\e^{2x+3}.
\]

Comme l'exponentielle est strictement croissante :
\[
x^2<2x+3.
\]

Donc :
\[
x^2-2x-3<0.
\]

On factorise :
\[
x^2-2x-3=(x-3)(x+1).
\]

On étudie le signe :
\[
(x-3)(x+1)<0
\]
pour :
\[
-1<x<3.
\]

Donc :
\[
S=]-1;3[.
\]
\end{exemplebox}

\subsection*{C. Changement d'inconnue simple}

Certaines équations peuvent être résolues en posant :
\[
X=\e^x.
\]

Comme $\e^x>0$, on a toujours :
\[
X>0.
\]

\begin{exemplebox}
Résoudre :
\[
\e^{2x}-3\e^x+2=0.
\]

On remarque que :
\[
\e^{2x}=(\e^x)^2.
\]

On pose :
\[
X=\e^x.
\]

L'équation devient :
\[
X^2-3X+2=0.
\]

On factorise :
\[
X^2-3X+2=(X-1)(X-2).
\]

Donc :
\[
X=1
\quad \text{ou} \quad
X=2.
\]

Si $X=1$, alors :
\[
\e^x=1,
\]
donc :
\[
x=0.
\]

Si $X=2$, l'équation :
\[
\e^x=2
\]
n'a pas de solution exprimable avec les outils de Première sans logarithme. On peut seulement dire qu'elle admet une unique solution réelle, car la fonction exponentielle est strictement croissante et prend toutes les valeurs positives.

Au niveau Première, on retient donc :
\[
x=0
\]
comme solution exacte accessible, et on peut approcher l'autre solution à la calculatrice.
\end{exemplebox}

%================================================
\section{Étude complète de fonctions avec exponentielle}

\subsection*{A. Méthode générale}

\begin{methode}
Pour étudier une fonction contenant une exponentielle :
\begin{enumerate}
    \item Déterminer l'ensemble de définition.
    \item Calculer la dérivée.
    \item Factoriser la dérivée.
    \item Utiliser $\e^x>0$ pour simplifier l'étude du signe.
    \item Construire le tableau de variations.
    \item Calculer les valeurs utiles : extremums, images, intersections.
    \item Rédiger une conclusion claire.
\end{enumerate}
\end{methode}

\subsection*{B. Exemple complet}

\begin{exemplebox}
On considère la fonction définie sur $\R$ par :
\[
f(x)=(x-2)\e^x.
\]

\textbf{1. Domaine de définition}

La fonction $x\mapsto x-2$ est définie sur $\R$ et la fonction exponentielle aussi. Donc $f$ est définie sur $\R$.

\textbf{2. Dérivée}

On pose :
\[
u(x)=x-2,
\qquad
v(x)=\e^x.
\]

Alors :
\[
u'(x)=1,
\qquad
v'(x)=\e^x.
\]

Donc :
\[
f'(x)=u'(x)v(x)+u(x)v'(x).
\]

Ainsi :
\[
f'(x)=\e^x+(x-2)\e^x.
\]

On factorise :
\[
f'(x)=(x-1)\e^x.
\]

\textbf{3. Signe de la dérivée}

Comme :
\[
\e^x>0,
\]
le signe de $f'(x)$ est celui de $x-1$.

Donc :
\[
f'(x)<0 \text{ si } x<1,
\]
\[
f'(1)=0,
\]
\[
f'(x)>0 \text{ si } x>1.
\]

\textbf{4. Variations}

La fonction $f$ est décroissante sur $]-\infty;1]$ puis croissante sur $[1;+\infty[$.

\textbf{5. Minimum}

On calcule :
\[
f(1)=(1-2)\e^1=-\e.
\]

Donc $f$ admet un minimum égal à $-\e$ en $x=1$.
\end{exemplebox}

%================================================
\section{Exponentielle et suites géométriques}

\subsection*{A. Rappel sur les suites géométriques}

\begin{definitionbox}
Une suite $(u_n)$ est géométrique lorsqu'il existe un réel $q$ tel que :
\[
u_{n+1}=qu_n
\]
pour tout entier naturel $n$.

Le nombre $q$ est appelé la raison de la suite.
\end{definitionbox}

Si $u_0$ est donné, alors :
\[
u_n=u_0q^n.
\]

\subsection*{B. Suites contenant l'exponentielle}

On peut rencontrer des suites du type :
\[
u_n=\e^{an+b}.
\]

\begin{exemplebox}
Soit :
\[
u_n=\e^{2n-1}.
\]

On étudie le quotient :
\[
\frac{u_{n+1}}{u_n}
=
\frac{\e^{2(n+1)-1}}{\e^{2n-1}}.
\]

On simplifie :
\[
\frac{u_{n+1}}{u_n}
=
\frac{\e^{2n+1}}{\e^{2n-1}}
=
\e^{(2n+1)-(2n-1)}
=
\e^2.
\]

Donc :
\[
u_{n+1}=\e^2u_n.
\]

La suite est géométrique de raison $\e^2$.
\end{exemplebox}

\subsection*{C. Sens de variation}

\begin{propriete}
La suite :
\[
u_n=\e^{an+b}
\]
est :
\begin{itemize}
    \item croissante si $a>0$ ;
    \item constante si $a=0$ ;
    \item décroissante si $a<0$.
\end{itemize}
\end{propriete}

\begin{proof}
On sait que la fonction exponentielle est strictement croissante.

Si $a>0$, alors $an+b$ augmente lorsque $n$ augmente, donc $\e^{an+b}$ augmente.

Si $a<0$, alors $an+b$ diminue lorsque $n$ augmente, donc $\e^{an+b}$ diminue.
\end{proof}

%================================================
\section{Modélisation avec l'exponentielle}

\subsection*{A. Modèle général}

Une fonction exponentielle de modélisation est souvent de la forme :
\[
f(t)=C\e^{kt}.
\]

\begin{itemize}
    \item $C$ représente la quantité initiale ;
    \item $k$ représente le taux continu d'évolution ;
    \item si $k>0$, la fonction est croissante ;
    \item si $k<0$, la fonction est décroissante.
\end{itemize}

\begin{retenir}
Un modèle exponentiel est adapté lorsqu'une quantité varie proportionnellement à la quantité présente.
\end{retenir}

\subsection*{B. Exemple de croissance}

\begin{exemplebox}
Une population est modélisée par :
\[
P(t)=1200\e^{0,03t},
\]
où $t$ est exprimé en années.

La population initiale est :
\[
P(0)=1200\e^0=1200.
\]

Comme $0,03>0$, la population augmente.

Au bout de $10$ ans :
\[
P(10)=1200\e^{0,3}.
\]

À la calculatrice :
\[
P(10)\approx 1200\times 1,35=1620.
\]

La population est donc d'environ $1620$ individus.
\end{exemplebox}

\subsection*{C. Exemple de décroissance}

\begin{exemplebox}
Une substance radioactive est modélisée par :
\[
Q(t)=500\e^{-0,2t}.
\]

La quantité initiale est :
\[
Q(0)=500.
\]

Comme $-0,2<0$, la quantité diminue.

Au bout de $5$ heures :
\[
Q(5)=500\e^{-1}.
\]

Or :
\[
\e^{-1}=\frac{1}{\e}\approx 0,368.
\]

Donc :
\[
Q(5)\approx 500\times 0,368=184.
\]

Il reste environ $184$ unités de substance.
\end{exemplebox}

%================================================
\section{Méthodes générales à maîtriser}

\subsection*{A. Simplifier une expression exponentielle}

\begin{methode}
\[
\e^a\e^b=\e^{a+b},
\qquad
\frac{\e^a}{\e^b}=\e^{a-b},
\qquad
(\e^a)^n=\e^{na}.
\]
\end{methode}

\subsection*{B. Comparer deux exponentielles}

\begin{methode}
Comme la fonction exponentielle est strictement croissante :
\[
\e^A<\e^B \Longleftrightarrow A<B.
\]

On compare donc les exposants.
\end{methode}

\subsection*{C. Étudier une dérivée factorisée}

\begin{methode}
Si :
\[
f'(x)=P(x)\e^x,
\]
alors comme :
\[
\e^x>0,
\]
le signe de $f'(x)$ est celui de $P(x)$.
\end{methode}

\subsection*{D. Reconnaître un modèle exponentiel}

\begin{methode}
Une situation peut être modélisée par une exponentielle lorsque :
\begin{itemize}
    \item la quantité évolue en pourcentage ;
    \item la variation est proportionnelle à la quantité présente ;
    \item il y a une croissance ou une décroissance multiplicative ;
    \item le modèle proposé est de la forme $C\e^{kt}$.
\end{itemize}
\end{methode}

%================================================
\section{Exercices progressifs}

\begin{exercice}[Calculs simples]
Simplifier les expressions suivantes :
\begin{enumerate}[label=\alph*)]
    \item $\e^2\e^5$
    \item $\dfrac{\e^7}{\e^3}$
    \item $\e^{x+1}\e^{2x-4}$
    \item $\dfrac{\e^{3x+2}}{\e^{x-1}}$
    \item $\e^{-x}\e^{4x+3}$
\end{enumerate}
\end{exercice}

\begin{exercice}[Équations]
Résoudre dans $\R$ :
\begin{enumerate}[label=\alph*)]
    \item $\e^x=\e^4$
    \item $\e^{2x-1}=\e^7$
    \item $\e^{x^2}=\e^{3x+4}$
    \item $\e^{x^2-5x}=\e^{-6}$
\end{enumerate}
\end{exercice}

\begin{exercice}[Inéquations]
Résoudre dans $\R$ :
\begin{enumerate}[label=\alph*)]
    \item $\e^x<\e^3$
    \item $\e^{2x+1}\geq \e^5$
    \item $\e^{x^2}<\e^{4}$
    \item $\e^{x^2-3x}\leq \e^{0}$
\end{enumerate}
\end{exercice}

\begin{exercice}[Dérivation]
Dériver les fonctions suivantes :
\begin{enumerate}[label=\alph*)]
    \item $f(x)=4\e^x-7$
    \item $g(x)=\e^x+3x-2$
    \item $h(x)=x\e^x$
    \item $p(x)=(x^2+1)\e^x$
\end{enumerate}
\end{exercice}

\begin{exercice}[Étude de variations]
On considère la fonction :
\[
f(x)=(x+3)\e^x.
\]
\begin{enumerate}
    \item Déterminer l'ensemble de définition de $f$.
    \item Calculer $f'(x)$.
    \item Factoriser $f'(x)$.
    \item Étudier le signe de $f'(x)$.
    \item Dresser le tableau de variations de $f$.
    \item Déterminer l'éventuel extremum de $f$.
\end{enumerate}
\end{exercice}

\begin{exercice}[Modélisation]
Une population de bactéries est modélisée par :
\[
N(t)=800\e^{0,4t},
\]
où $t$ est exprimé en heures.
\begin{enumerate}
    \item Donner la population initiale.
    \item Dire si la population augmente ou diminue.
    \item Calculer une valeur approchée de la population au bout de $3$ heures.
    \item Interpréter le coefficient $0,4$.
\end{enumerate}
\end{exercice}

%================================================
\section{Grand devoir bilan}

\subsection*{Exercice 1 : Calculs et équations}

\begin{enumerate}
    \item Simplifier :
    \[
    A=\frac{\e^{2x+3}\e^{x-1}}{\e^{4x}}.
    \]
    \item Résoudre :
    \[
    \e^{3x-2}=\e^{x+6}.
    \]
    \item Résoudre :
    \[
    \e^{x^2-2x}<\e^3.
    \]
\end{enumerate}

\subsection*{Exercice 2 : Étude de fonction}

On considère la fonction définie sur $\R$ par :
\[
f(x)=(x-1)\e^x.
\]

\begin{enumerate}
    \item Calculer $f'(x)$.
    \item Montrer que :
    \[
    f'(x)=x\e^x.
    \]
    \item Étudier le signe de $f'(x)$.
    \item Dresser le tableau de variations de $f$.
    \item Déterminer le minimum de $f$.
\end{enumerate}

\subsection*{Exercice 3 : Modélisation}

Une quantité de médicament présente dans le sang est modélisée par :
\[
M(t)=60\e^{-0,25t},
\]
où $t$ est exprimé en heures.

\begin{enumerate}
    \item Calculer $M(0)$.
    \item Justifier que la quantité diminue.
    \item Calculer une valeur approchée de $M(4)$.
    \item Interpréter le résultat dans le contexte.
\end{enumerate}

%================================================
\section{Corrections détaillées}

\begin{correction}
\begin{enumerate}[label=\alph*)]
    \item $\e^2\e^5=\e^{2+5}=\e^7$.
    \item $\dfrac{\e^7}{\e^3}=\e^{7-3}=\e^4$.
    \item $\e^{x+1}\e^{2x-4}=\e^{(x+1)+(2x-4)}=\e^{3x-3}$.
    \item $\dfrac{\e^{3x+2}}{\e^{x-1}}=\e^{(3x+2)-(x-1)}=\e^{2x+3}$.
    \item $\e^{-x}\e^{4x+3}=\e^{-x+4x+3}=\e^{3x+3}$.
\end{enumerate}
\end{correction}

\begin{correction}
\begin{enumerate}[label=\alph*)]
    \item $\e^x=\e^4$ donc $x=4$.
    \item $\e^{2x-1}=\e^7$ donc $2x-1=7$, d'où $x=4$.
    \item $\e^{x^2}=\e^{3x+4}$ donc $x^2=3x+4$.
    Ainsi :
    \[
    x^2-3x-4=0.
    \]
    On factorise :
    \[
    (x-4)(x+1)=0.
    \]
    Donc :
    \[
    x=4 \quad \text{ou} \quad x=-1.
    \]
    \item $\e^{x^2-5x}=\e^{-6}$ donc :
    \[
    x^2-5x=-6.
    \]
    Ainsi :
    \[
    x^2-5x+6=0.
    \]
    On factorise :
    \[
    (x-2)(x-3)=0.
    \]
    Donc :
    \[
    x=2 \quad \text{ou} \quad x=3.
    \]
\end{enumerate}
\end{correction}

\begin{correction}
\begin{enumerate}[label=\alph*)]
    \item $\e^x<\e^3$ équivaut à $x<3$.
    Donc :
    \[
    S=]-\infty;3[.
    \]
    \item $\e^{2x+1}\geq \e^5$ équivaut à :
    \[
    2x+1\geq 5.
    \]
    Donc :
    \[
    2x\geq 4,
    \]
    puis :
    \[
    x\geq 2.
    \]
    \item $\e^{x^2}<\e^4$ équivaut à :
    \[
    x^2<4.
    \]
    Donc :
    \[
    -2<x<2.
    \]
    \item $\e^{x^2-3x}\leq \e^0$ équivaut à :
    \[
    x^2-3x\leq 0.
    \]
    On factorise :
    \[
    x(x-3)\leq 0.
    \]
    Donc :
    \[
    S=[0;3].
    \]
\end{enumerate}
\end{correction}

\begin{correction}
\begin{enumerate}[label=\alph*)]
    \item $f(x)=4\e^x-7$, donc :
    \[
    f'(x)=4\e^x.
    \]
    \item $g(x)=\e^x+3x-2$, donc :
    \[
    g'(x)=\e^x+3.
    \]
    \item $h(x)=x\e^x$.
    Par produit :
    \[
    h'(x)=1\times \e^x+x\times \e^x=(x+1)\e^x.
    \]
    \item $p(x)=(x^2+1)\e^x$.
    Par produit :
    \[
    p'(x)=2x\e^x+(x^2+1)\e^x.
    \]
    Donc :
    \[
    p'(x)=(x^2+2x+1)\e^x=(x+1)^2\e^x.
    \]
\end{enumerate}
\end{correction}

\begin{correction}
On considère :
\[
f(x)=(x+3)\e^x.
\]

\begin{enumerate}
    \item La fonction est définie sur $\R$.
    \item On dérive par produit :
    \[
    f'(x)=1\times \e^x+(x+3)\e^x.
    \]
    \item Donc :
    \[
    f'(x)=(x+4)\e^x.
    \]
    \item Comme $\e^x>0$, le signe de $f'(x)$ est celui de $x+4$.
    Ainsi :
    \[
    f'(x)<0 \text{ si } x<-4,
    \]
    \[
    f'(x)=0 \text{ si } x=-4,
    \]
    \[
    f'(x)>0 \text{ si } x>-4.
    \]
    \item La fonction est décroissante sur $]-\infty;-4]$, puis croissante sur $[-4;+\infty[$.
    \item Elle admet un minimum en $x=-4$ :
    \[
    f(-4)=(-4+3)\e^{-4}=-\e^{-4}.
    \]
\end{enumerate}
\end{correction}

\begin{correction}
On a :
\[
N(t)=800\e^{0,4t}.
\]

\begin{enumerate}
    \item La population initiale est :
    \[
    N(0)=800\e^0=800.
    \]
    \item Comme $0,4>0$, la fonction modélise une croissance.
    \item Au bout de $3$ heures :
    \[
    N(3)=800\e^{1,2}.
    \]
    À la calculatrice :
    \[
    \e^{1,2}\approx 3,32.
    \]
    Donc :
    \[
    N(3)\approx 800\times 3,32=2656.
    \]
    \item Le coefficient $0,4$ mesure la rapidité de croissance continue de la population.
\end{enumerate}
\end{correction}

%================================================
\section{Correction du devoir bilan}

\subsection*{Correction de l'exercice 1}

\begin{enumerate}
    \item
    \[
    A=\frac{\e^{2x+3}\e^{x-1}}{\e^{4x}}.
    \]

    Au numérateur :
    \[
    \e^{2x+3}\e^{x-1}=\e^{3x+2}.
    \]

    Donc :
    \[
    A=\frac{\e^{3x+2}}{\e^{4x}}=\e^{3x+2-4x}.
    \]

    Ainsi :
    \[
    A=\e^{2-x}.
    \]

    \item
    \[
    \e^{3x-2}=\e^{x+6}.
    \]

    Par injectivité de l'exponentielle :
    \[
    3x-2=x+6.
    \]

    Donc :
    \[
    2x=8,
    \]
    d'où :
    \[
    x=4.
    \]

    \item
    \[
    \e^{x^2-2x}<\e^3
    \]
    équivaut à :
    \[
    x^2-2x<3.
    \]

    Donc :
    \[
    x^2-2x-3<0.
    \]

    On factorise :
    \[
    x^2-2x-3=(x-3)(x+1).
    \]

    Le produit est négatif entre ses deux racines :
    \[
    -1<x<3.
    \]

    Donc :
    \[
    S=]-1;3[.
    \]
\end{enumerate}

\subsection*{Correction de l'exercice 2}

On considère :
\[
f(x)=(x-1)\e^x.
\]

\begin{enumerate}
    \item On dérive par produit :
    \[
    f'(x)=1\times \e^x+(x-1)\e^x.
    \]

    \item On factorise :
    \[
    f'(x)=\left(1+x-1\right)\e^x=x\e^x.
    \]

    \item Comme $\e^x>0$, le signe de $f'(x)$ est celui de $x$.

    Donc :
    \[
    f'(x)<0 \text{ si } x<0,
    \]
    \[
    f'(0)=0,
    \]
    \[
    f'(x)>0 \text{ si } x>0.
    \]

    \item La fonction $f$ est décroissante sur $]-\infty;0]$ puis croissante sur $[0;+\infty[$.

    \item Le minimum est atteint en $x=0$ :
    \[
    f(0)=(0-1)\e^0=-1.
    \]

    Donc $f$ admet un minimum égal à $-1$.
\end{enumerate}

\subsection*{Correction de l'exercice 3}

On a :
\[
M(t)=60\e^{-0,25t}.
\]

\begin{enumerate}
    \item
    \[
    M(0)=60\e^0=60.
    \]

    La quantité initiale est donc de $60$ unités.

    \item Le coefficient de $t$ dans l'exposant est :
    \[
    -0,25.
    \]

    Il est négatif, donc la quantité diminue.

    \item
    \[
    M(4)=60\e^{-0,25\times 4}=60\e^{-1}.
    \]

    Or :
    \[
    \e^{-1}\approx 0,368.
    \]

    Donc :
    \[
    M(4)\approx 60\times 0,368=22,08.
    \]

    \item Après $4$ heures, il reste environ $22,1$ unités de médicament dans le sang.
\end{enumerate}

%================================================
\section{Fiche de synthèse finale}

\begin{retenir}
\textbf{Définition :}
\[
\exp'(x)=\exp(x),
\qquad
\exp(0)=1.
\]

\textbf{Notation :}
\[
\exp(x)=\e^x.
\]

\textbf{Valeurs importantes :}
\[
\e^0=1,
\qquad
\e^1=\e,
\qquad
\e^{-1}=\frac{1}{\e}.
\]

\textbf{Propriétés de calcul :}
\[
\e^{a+b}=\e^a\e^b,
\]
\[
\e^{-a}=\frac{1}{\e^a},
\]
\[
\e^{a-b}=\frac{\e^a}{\e^b}.
\]

\textbf{Signe :}
\[
\e^x>0.
\]

\textbf{Dérivée :}
\[
(\e^x)'=\e^x.
\]

\textbf{Variations :}
La fonction exponentielle est strictement croissante sur $\R$.

\textbf{Injectivité :}
\[
\e^A=\e^B \Longleftrightarrow A=B.
\]

\textbf{Inéquations :}
\[
\e^A<\e^B \Longleftrightarrow A<B.
\]
\end{retenir}

%================================================
\section{Erreurs les plus fréquentes}

\begin{erreur}
Confondre :
\[
\e^{a+b}
\]
avec :
\[
\e^a+\e^b.
\]

La bonne formule est :
\[
\e^{a+b}=\e^a\e^b.
\]
\end{erreur}

\begin{erreur}
Oublier que :
\[
\e^x>0
\]
pour tout réel $x$.

Dans une dérivée comme :
\[
f'(x)=(x-3)\e^x,
\]
le signe dépend seulement de $x-3$.
\end{erreur}

\begin{erreur}
Écrire que :
\[
\e^x=0
\]
a une solution.

C'est faux : l'exponentielle ne s'annule jamais.
\end{erreur}

\begin{erreur}
Utiliser le logarithme alors qu'il n'a pas encore été introduit dans le cours.

En Première, on résout surtout les équations de la forme :
\[
\e^A=\e^B.
\]
\end{erreur}

%================================================
\section{Ouverture vers la Terminale}

En Terminale, l'étude de l'exponentielle sera approfondie avec :
\begin{itemize}
    \item les limites de la fonction exponentielle ;
    \item la comparaison entre exponentielle et puissances ;
    \item la fonction logarithme népérien, réciproque de l'exponentielle ;
    \item la convexité ;
    \item les équations différentielles du type $y'=ay$ ;
    \item des modèles plus avancés de croissance et de décroissance.
\end{itemize}

\begin{pointbac}
La maîtrise de l'exponentielle en Première est essentielle pour réussir les études de fonctions en Terminale.
\end{pointbac}

\end{document}